\section{Dataset and evaluation protocol}\label{sec:protocol}

\subsection{Splits and subsets}
Table~\ref{tab:splits} summarises the splits. All development decisions used TRAIN and VAL only.

\begin{table}[t]
\centering
\caption{SeaClips splits as used in this study. Frame counts were taken from the archive file list.}
\label{tab:splits}
\begin{tabular}{lrrl}
\toprule
Split & Videos & Frames & Use \\
\midrule
TRAIN & 47 & 20{,}310 & training \\
VAL & 9 & 4{,}266 & two 432/426-frame subsets for gates \\
TEST & 18 & 7{,}030 & single final evaluation \\
\bottomrule
\end{tabular}
\end{table}

Running each gate on all VAL frames was too slow for tiled inference on one GPU (about 2\,s per frame), so we defined two disjoint subsets of our own. Within each video, frames are sorted by frame index, and those with index $\equiv 0 \pmod{10}$ (432 frames, the \emph{evaluation} subset) or $\equiv 5 \pmod{10}$ (426 frames, the \emph{selection and replication} subset) are kept. On the R50 baseline, the evaluation subset reproduces full-VAL \mapthree{} to within 0.0004. Because frames are recorded at 30\,frames per second and are strongly correlated, intervals are computed by bootstrap over the nine VAL videos rather than over frames.

\subsection{Metric}
We use pycocotools \citep{lin2014coco} with an IoU threshold of 0.30, at most 100 detections per image, and boxes in original image coordinates, as specified by the benchmark. \mapthree{} is the mean of per-class AP over the classes that have ground truth: seven on TEST and six on VAL, which contains no LEISURE\_VEHICLE. We also report per-class AP, AP for COCO small ($<32^2$\,px), medium and large objects, class-agnostic mAP (mAPca@0.3) and mAP@0.5.

\subsection{Pre-registration and TEST lock}
Each gate was described in a pre-declaration file before it ran: the hypothesis, the data, the arms, a pass rule with numeric thresholds, and the consequence of each outcome. The pre-declaration recorded the SHA-256 of the script that would run. Every long run was preceded by a self-test on synthetic data and a short GPU smoke run that checked resumption, memory and throughput.

Before any final model existed, the TEST label file was hashed without being parsed. The evaluation tool refused any other file, refused to score until every model had been exported, and wrote a lock file after its single run.

\subsection{Statistics}
As in the benchmark paper, we compare the mean of three seeds against each published mean $\pm$ standard deviation ($n=3$) with a one-tailed Welch t-test \citep{welch1947} computed from summary statistics. Degrees of freedom follow Welch--Satterthwaite, and a Bonferroni correction is applied over three comparisons ($\alpha = 0.05/3$). Before TEST was opened we computed that, given the published standard deviation of TransVOD++, a significant difference would require a three-seed mean near 40 \mapthree{}; this fixed the wording of any claim in advance.
